% Drawing for measurement `jfet-runs`
% Stands for the deck lines:
%   VD d 0 10
%   VG g 0 0
%   J1 d g 0 DUT
%   .op
\documentclass[border=8pt]{standalone}
\usepackage[siunitx,RPvoltages,european]{circuitikz}
\ctikzset{bipoles/length=0.9cm, label/align=straight, inductors/scale=0.8}
\ctikzset{inductor=american}
\usepackage{amsmath}
\begin{document}
\begin{circuitikz}[font=\footnotesize, line width=0.5pt]
\draw (0,0) node[njfet](Q){};
\node[anchor=west] at ($(Q.D)!0.5!(Q.S) + (0.55,0)$) {DUT};
\draw (Q.S) to[short] ++(0,-0.5) node[ground]{};
\draw (Q.D) to[short] ++(0,0.3) to[rmeter, t=A, l={$I_D$}] ++(0,1.4) to[short] ++(2.0,0) to[vsource, l={$V_{DS}=\SI{10}{V}$}, invert] ++(0,-2.4) node[ground]{};
\draw (Q.G) to[short] ++(-0.5,0) to[rmeter, t=A, l={$I_G$}] ++(-1.4,0) to[short] ++(-0.5,0) coordinate(gl) to[vsource, l_={$V_{GS}=0$}, invert] ++(0,-1.9) node[ground]{};
\end{circuitikz}
\end{document}
